% Irodov Part III §3.3 122--130; PDF pp.117--119,
% printed pp.108--110. Fig.3.24--3.29 checked.
\electrocapproblem{122}{108}
图3.24：左、右两个电源电动势都为
$\mathcal E$，电容器 $C_1$ 接在右上分支，
$C_2$ 接在中间竖支路。原先开关 Sw 断开，
闭合后求沿图中箭头方向流过截面1、2的电荷量。
\exsolution
取底部左侧公共点为零电势。
闭合前，左侧电源未接入，
$C_1,C_2$ 串联接右侧电源，
初始串联电荷
$Q_0=\mathcal E C_1C_2/(C_1+C_2)$。
闭合后左电源将中间上节点钳在
$+\mathcal E$，右电源上端同样为
$+\mathcal E$，故 $C_1$ 两端电压归零，
$C_2$ 两端电压为 $\mathcal E$。
截面1按左箭头的累计电荷
等于 $C_2$ 的最终电荷：
\[
\boxed{q_1=\mathcal E C_2}.
\]
截面2按图中右箭头，$C_1$ 的
初有电荷 $Q_0$ 被消去，对该方向贡献为负：
\[
\boxed{q_2=-\frac{\mathcal E C_1C_2}{C_1+C_2}}.
\]
\excheck 截面电荷是开关动作前后的\emph{累计通过量}，
不等于最后流动的稳态电流（稳态电流为零）。

\electrocapproblem{123}{108}
图3.25：左侧两电源均为
$\mathcal E=60\,\mathrm V$，右侧两只电容器
$C_1=2.0\,\mathrm{\mu F}$、
$C_2=3.0\,\mathrm{\mu F}$。
闭合中间 Sw 后，求沿箭头1、2、3
的累计电荷量。
\exsolution
闭合前，右侧串联电容器跨越两个电源
给出的总电压 $2\mathcal E$，
故每只电容器的电荷模
$Q_0=2\mathcal E C_1C_2/(C_1+C_2)$。
闭合后两只电容器各跨一只电动势
$\mathcal E$，分别带电
$Q_{1f}=\mathcal E C_1$、
$Q_{2f}=\mathcal E C_2$。
由各导线截面的电荷连续性，依原书箭头：
\[
\boxed{q_1=\mathcal E C_1-Q_0
=\mathcal E C_1\frac{C_1-C_2}{C_1+C_2}},
\]
\[
\boxed{q_2=-\mathcal E C_2
\frac{C_2-C_1}{C_1+C_2}
=\mathcal E C_2\frac{C_1-C_2}{C_1+C_2}},
\]
\[
\boxed{q_3=\mathcal E(C_2-C_1)}.
\]
代入数据：
\[
\boxed{q_1=-24\,\mathrm{\mu C},\qquad
q_2=-36\,\mathrm{\mu C},\qquad
q_3=+60\,\mathrm{\mu C}}.
\]
\excheck 两个负号说明实际电荷流过截面1、2时
与图中所指箭头相反；第3截面为正。

\electrocapproblem{124}{108--109}
原书图3.26：上支路为
$C_1,C_2$ 两只串联电容器，
中点 A 通过电容器 $C_3$ 连接下方 B，
底部分别有电动势 $\mathcal E_1,\mathcal E_2$。
求 $\varphi_A-\varphi_B$。
\exsolution
取 $\varphi_B=0$。
按图中电源极性，
$C_1$ 左端电势为 $-\mathcal E_1$，
$C_2$ 右端电势为 $+\mathcal E_2$。
A 是孤立浮置金属节点，初始无净自由电荷，
其三个电容器极板电荷代数和为零：
\[
C_1(\varphi_A+\mathcal E_1)
+C_2(\varphi_A-\mathcal E_2)
+C_3\varphi_A=0.
\]
于是
\[
\boxed{\varphi_A-\varphi_B=
\frac{C_2\mathcal E_2-C_1\mathcal E_1}
{C_1+C_2+C_3}}.
\]
\excheck $C_3$ 越大，中间节点相对于 B
的电势差越小，符合把 A 拉向 B 电势的直觉。

\electrocapproblem{125}{109}
原书图3.27：三个电源
$\mathcal E_1,\mathcal E_2,\mathcal E_3$
沿一竖支路串接，三只电容器
$C_1,C_2,C_3$ 分别把各层节点
1、2、3 接到公共点 O。
规定 $\varphi_O=0$，求
$\varphi_1,\varphi_2,\varphi_3$。
\exsolution
按图中各电动势极性设定，
三个左侧节点之间满足
$\varphi_i=\Phi-\mathcal E_i$，
其中 $\Phi$ 为共同参考偏移。
整个左侧电源组与外界通过电容耦合，
总净电荷守恒：
\[
C_1\varphi_1+C_2\varphi_2+C_3\varphi_3=0.
\]
故
$\Phi=(C_1\mathcal E_1+C_2\mathcal E_2+
C_3\mathcal E_3)/(C_1+C_2+C_3)$。
定义 $C_\Sigma=C_1+C_2+C_3$，得到
\[
\boxed{\varphi_1=
\frac{C_2\mathcal E_2+C_3\mathcal E_3-
(C_2+C_3)\mathcal E_1}{C_\Sigma}},
\]
\[
\boxed{\varphi_2=
\frac{C_1\mathcal E_1+C_3\mathcal E_3-
(C_1+C_3)\mathcal E_2}{C_\Sigma}},
\]
\[
\boxed{\varphi_3=
\frac{C_1\mathcal E_1+C_2\mathcal E_2-
(C_1+C_2)\mathcal E_3}{C_\Sigma}}.
\]
\excheck 将三式分别乘 $C_1,C_2,C_3$ 再相加，
应得零，满足浮置导体组的总电荷守恒。

\electrocapproblem{126}{109}
原书图3.28电容桥：
A—B 之间上支路依次为 $C_1,C_2$，
下支路依次为 $C_2,C_1$；
两串联支路中点通过一只 $C_3$ 桥接。
求总等效电容。
\exsolution
设 A 电势为 $V$，B 为零。
上下支路的中点电势分别为 $x,y$。
电荷守恒：
\[
C_1(x-V)+C_2x+C_3(x-y)=0,
\]
\[
C_2(y-V)+C_1y+C_3(y-x)=0.
\]
两式相加得 $x+y=V$；
相减得到
\[
x-y=\frac{(C_1-C_2)V}{C_1+C_2+2C_3}.
\]
A 端净电荷为
$Q=C_1(V-x)+C_2(V-y)$，
代入得
\[
C_{AB}=\frac QV
=\frac{2C_1C_2+C_3(C_1+C_2)}
{C_1+C_2+2C_3}.
\]
因此
\[
\boxed{C_{AB}=
\frac{2C_1C_2+C_3(C_1+C_2)}
{C_1+C_2+2C_3}}.
\]
\excheck $C_3\to0$ 时恢复
两条串联支路并联
$2C_1C_2/(C_1+C_2)$；
$C_1=C_2$ 时桥中点同势，结果为 $C_1$，
不依赖 $C_3$。

\electrocapproblem{127}{109}
图3.29(a)、(b)、(c)：
边长 $a$ 的正方形四顶点分别放置
大小均为 $q$ 的点电荷；(a) 全为正，
(b) 相邻顶点正负交错，
(c) 上两顶点正、下两顶点负。
求每种构型四电荷的相互作用总能。
\exsolution
离散点电荷的相互作用能
$W=k\sum_{i<j}q_iq_j/r_{ij}$，
正方形有4条边（距离 $a$）
及2条对角线（距离 $\sqrt2a$）。
(a) 全部相同，4条边贡献
$4kq^2/a$、对角线贡献
$\sqrt2kq^2/a$，所以
\[
\boxed{W_a=(4+\sqrt2)\frac{q^2}
{4\pi\varepsilon_0a}}.
\]
(b) 四条边均为异号，贡献
$-4kq^2/a$；两条对角线同号，
贡献 $+\sqrt2kq^2/a$：
\[
\boxed{W_b=(\sqrt2-4)\frac{q^2}
{4\pi\varepsilon_0a}}.
\]
(c) 两条水平边同号、两条竖直边异号，
边贡献抵消；两条对角线异号：
\[
\boxed{W_c=-\sqrt2\frac{q^2}
{4\pi\varepsilon_0a}}.
\]
\excheck 相互作用能只对每一\emph{无序电荷对}
计数一次，不加点电荷本身的无穷自能。

\electrocapproblem{128}{109}
无限直线链由 $q,-q,q,-q,\ldots$
交替排列，相邻距离为 $a$。
求链中任意一个电荷与其余电荷的相互作用能。
\exsolution
取其中一个 $+q$ 位于原点。
距离原点左右两侧各有
$n a$ 处电荷 $(-1)^nq$，$n\ge1$。
原点电荷的\emph{与其余所有电荷作用能}
为 $q$ 乘外部电势：
\[
W=2k\frac{q^2}{a}
\sum_{n=1}^\infty\frac{(-1)^n}{n}.
\]
利用
$\ln(1+x)=\sum_{n\ge1}(-1)^{n-1}x^n/n$
在 $x=1$ 的极限，
$\sum_{n\ge1}(-1)^n/n=-\ln2$，
得
\[
\boxed{W=-\frac{2q^2\ln2}
{4\pi\varepsilon_0a}}.
\]
\excheck 本题问\emph{指定一个电荷与其余电荷}
的能量。如果要求无限链的平均总能
按每个电荷分摊，应再除以2以避免成对重复计数。

\electrocapproblem{129}{109}
点电荷 $q$ 距无限接地金属导体平面为
$\ell$。求该电荷与导体表面
\emph{感应电荷的相互作用能}。
\exsolution
导体的感应电荷在真电荷位置所产生的电势
等于镜像电荷 $-q$ 在距离 $2\ell$ 处的电势：
\[
\varphi_{\rm ind}=
-\frac{q}{4\pi\varepsilon_0(2\ell)}
=-\frac{q}{8\pi\varepsilon_0\ell}.
\]
本题直接求\emph{真电荷与给定感应分布}
之间的双体相互作用项
$W_{\rm int}=q\varphi_{\rm ind}$，所以
\[
\boxed{W_{\rm int}=
-\frac{q^2}{8\pi\varepsilon_0\ell}}.
\]
\excheck 这与从无限远\emph{组装}
真电荷—导体完整体系所需的
$-\frac{q^2}{16\pi\varepsilon_0\ell}$
不相同：后者包含诱导分布随运动改变的
$1/2$ 能量因子。原书本题明确问相互作用项。

\electrocapproblem{130}{109--110}
两颗球形对称带电小球，净电荷为
$q_1,q_2$，两球中心距离为
$\ell$，彼此不重叠。
求两球间相互作用能。
\exsolution
对球外点，球对称电荷分布产生的静电势
与中心点电荷相同：
$\varphi_1(r)=kq_1/r$。
如果第二球完整位于第一球之外，
第一球的势在第二球的各球壳上按球对称平均，
等于其中心位置的势
$kq_1/\ell$（平均值定理）。
因此
\[
W_{12}=\int_{\text{球2}}\varphi_1(\bm r)
\,\mathrm dq_2(\bm r)
=q_2\frac{kq_1}{\ell},
\]
故
\[
\boxed{W_{12}=
\frac{q_1q_2}{4\pi\varepsilon_0\ell}}.
\]
\excheck 需两颗球的电荷分布均保持球对称；
若球是相互极化的导体，近距离时分布可能改变，
不能直接套此简式。
